% % methods.tex
% \begin{enumerate}[partopsep=0pt,itemsep=0pt,leftmargin=0.4cm,label=\protect\circled{\arabic*}]
%     \item Cada \textbf{individuo} representa una imagen compuesta por triángulos.
%           \begin{center}
%               \vspace{-1em}
%               \begin{tikzpicture}
%                   \node[] (imageA) at (0,0){\includegraphics[width=0.5\textwidth]{monalisa/fittest_1gen_50pop.png}};
%                   \node[below=1pt of imageA](captionA){\scriptsize Some caption.};
%                   % \draw[step=1cm,blue,thin] (0,0) grid (6,5);
%               \end{tikzpicture}
%               \vspace{-1em}
%           \end{center}
%     \item Se evalúa cada imagen mediante la métrica CIE Lab $\Delta CIE1994$ para medir la
%           \textbf{similitud de color}.\\ Se usa la ecuación \eqref{eq:CIE1994}
%     \item El AG selecciona (torneo), cruza y muta las mejores soluciones para mejorar la aproximación.
% \end{enumerate}
\begin{enumerate}[partopsep=0pt,itemsep=0pt,leftmargin=0.4cm,label=\protect\circled{\arabic*}]
    \item Each \textbf{individual} represents an image composed of triangles.
          \begin{center}
              \vspace{-1em}
              \begin{tikzpicture}
                  \node[] (imageA) at (0,0){\includegraphics[width=0.4\textwidth]{monalisa/fittest_1gen_50pop.png}};
                  \node[below=1pt of imageA](captionA){\small Example of a first individual.};
                  % \draw[step=1cm,blue,thin] (0,0) grid (6,5);
              \end{tikzpicture}
              \vspace{-1em}
          \end{center}
    \item Each image is evaluated using the CIE Lab $\Delta CIE1994$ metric to measure
          \textbf{color similarity} as the mean value.\\ Equation \eqref{eq:CIE1994} is used in each pixel.
    \item The GA selects (tournament), crosses, and mutates the best solutions to improve the approximation.
    \item The process is repeated over several generations until an optimal solution is reached.
    \item A stop condition is applied when fitness is $< 0.01$ for 3 consecutive individuals.
\end{enumerate}

% Para evaluar las capacidades de \boxi[colorD]{tcolorbox}, se utilizaron las siguientes métricas y metodologías:

% \begin{enumerate}[partopsep=0pt,itemsep=0pt,leftmargin=0.4cm,label=\protect\circled{\arabic*}]
%     \item \textbf{Tiempo de Elaboración}:
%           \begin{itemize}[topsep=0pt, partopsep=0pt,itemsep=1pt,leftmargin=0.4cm]
%               \item Se midió el tiempo total requerido para crear un póster desde cero utilizando \boxi[colorD]{tcolorbox}.
%               \item Se registraron los tiempos de compilación.
%           \end{itemize}

%     \item \textbf{Facilidad de Uso}:
%           \begin{itemize}[topsep=0pt, partopsep=0pt,itemsep=1pt,leftmargin=0.4cm]
%               \item Se evaluó la curva de aprendizaje para usuarios de diferentes niveles.
%           \end{itemize}

%     \item \textbf{Flexibilidad de Diseño}:
%           \begin{itemize}[topsep=0pt, partopsep=0pt,itemsep=1pt,leftmargin=0.4cm]
%               \item Se exploraron las opciones de personalización disponibles en \boxi[colorD]{tcolorbox}.
%               \item Se verificó la capacidad de implementar diseños complejos.
%           \end{itemize}

%     \item \textbf{Calidad del Resultado Final}:
%           \begin{itemize}[topsep=0pt, partopsep=0pt,itemsep=1pt,leftmargin=0.4cm]
%               \item Un panel de expertos evaluó la calidad visual del póster generado.
%               \item Se compararon aspectos como la resolución, alineación, consistencia tipográfica.
%           \end{itemize}

%     \item \textbf{Compatibilidad y Dependencias}:
%           \begin{itemize}[topsep=0pt, partopsep=0pt,itemsep=1pt,leftmargin=0.4cm]
%               \item Se revisó la compatibilidad de \boxi[colorD]{tcolorbox} con otros paquetes comunes de \LaTeX.
%               \item Se identificaron las dependencias adicionales requeridas.
%           \end{itemize}
%     \item \textbf{Documentación y Soporte}:
%           \begin{itemize}[topsep=0pt, partopsep=0pt,itemsep=1pt,leftmargin=0.4cm]
%               \item Se evaluó la calidad de la documentación oficial de \boxi[colorD]{tcolorbox}.
%               \item Se revisó la disponibilidad de tutoriales, ejemplos y comunidades de soporte en línea.
%           \end{itemize}
% \end{enumerate}
% % \begin{tikzpicture}
% %     \tikzstyle{textLabel} = [align=center,text width=8.75cm,rounded corners,fill=colorA!20,draw=colorA,line width=0.5 mm,text=colorA,minimum height=1.25cm]
% %     \node[textLabel] (aText) at (0,0){Se utilizó \texttt{tikz} para crear gráficos y diagramas.};
% %     \node[textLabel,below=2mm of aText] (bText){Se empleó \texttt{tcolorbox} para\\ estructurar el contenido en cajas.};
% %     \node[textLabel,below=2mm of bText] (cText){Comparamos los resultados entre \\ \texttt{tikzposter} y \texttt{beamer}.};
% %     \node[textLabel,below=2mm of cText] (dText){Elaboramos el mismo poster con ambos paquetes.};
% % \end{tikzpicture}